\begin{textsample}{POS Dim 1 – se – Score 31.00 – se/se\_em\_14.txt}  \label{ex:f1_pos_117}
3.3.3 Química A Química é \textbf{uma} Ciência \textbf{que} estuda os elementos representantes da matéria e suas transformações.
A compreensão dessa Ciência, relacionada com o seu desenvolvimento histórico e científico, \textbf{permeia} o campo das ideias em \textbf{que} se interligam os contextos sociais, políticos, culturais, religiosos, filosóficos e tecnológicos.
O início da Química vem desde os milhares de anos a.C e está ligada ao desenvolvimento da \textbf{humanidade}, por exemplo, desde a pré-história o \textbf{homem} aprendeu a manipular e transformar madeiras, objetos de pedra, ossos, rochas e peles em seu benefício próprio, bem como a descoberta do fogo.
Com esta descoberta foi possível provocar novas transformações na alimentação entre outros ( CHASSOT, 1998 ).
Os saberes e práticas ligadas à transformação da matéria, esteve presente nas diversas civilizações.
Caminhando \textbf{adiante}, os povos mesopotâmicos utilizaram as ligas metálicas para melhoria das armas e utensílios domésticos.
Os Egípcios desenvolveram técnicas para manipulação de metais, fabricação de vidro, porcelana, papiro, corantes, medicamentos, perfumes, cerveja e embalsamento de corpos.
Sendo assim, não é possível falar do surgimento da Química sem nos referirmos à Alquimia.
Alquimia é \textbf{uma} prática antiga, \textbf{que} se \textbf{baseava} em dois princípios fundamentais : ( 1 ) transformar os metais inferiores em ouro, ( Pedra Filosofal ), \textbf{baseada} nas ideias \textbf{que} Aristóteles e Empédocles tinham sobre a matéria e ( 2 ) manipulação do elixir da longa vida, ( Elixir da Vida Eterna ).
Alquimia tem \textbf{um} caráter místico, a palavra alquimia vem do árabe e significa ( AL\_KHMY ), surgiu das ciências da Mesopotâmia, Egito, Árabe e Pérsia.
A alquimia combinava química, física, filosofia, arte, metalurgia, medicina e religião ( GREENBERG, 2009 ).
Os Alquimistas foram os pioneiros no desenvolvimento de diversas técnicas de laboratório \textbf{que} são utilizadas até \textbf{hoje} ( claro \textbf{que} de forma aprimorada ).
Entre essas invenções estão o banho-maria, a destilação e a sublimação.
Entre as substâncias descobertas estão o álcool, ácido sulfúrico e ácido nítrico.
Eles se dedicavam de modo herméticos, envolvia muito mistério e segredos.
Na época não era aceita como ciências, era \textbf{visto} como bruxaria. [... ] eles buscavam no elixir da longa vida o \textbf{que} \textbf{hoje} se busca por Meio de remédios : melhorar a qualidade de vida e até prolongá -la.
A busca de novos materiais para o fabrico de vestuário e para Construção de habitações se assemelha ao \textbf{que} faziam os alquimistas \textbf{que} com a evaporação dos líquidos e a recalcinação de \textbf{sólidos} procuravam melhorar a qualidade das substâncias.
As retortas, os crisóis, os alambiques de então estão nos \textbf{modernos} Laboratórios de \textbf{hoje}, sob a forma de sofisticada aparelhagem de Vidros especiais ( CHASSOT, 2003, p. 119 ).
A partir da Alquimia surgiu a Iatroquímica, representada principalmente pela figura de Phillipus Aureolus Theophrastus Bombast von Hohenheim, mais \textbf{conhecido} como Paracelso.
Ele era \textbf{um} filósofo, médico e alquimista suíço e \textbf{que} se dedicava a combater doenças por meio da criação de medicamentos à base de fontes minerais.
No entanto, somente no \textbf{século} XVII, na Europa, é \textbf{que} a ciência Química teve seu cenário de desenvolvimento de modo de produção capitalista, atendendo aos interesses econômicos da classe dominante.
Na época ocorria a expansão da indústria, do comércio, da navegação e das técnicas militares. \textbf{Um} dos membros da ROYAL SOCIETY, Robert Boyle, tornou -se firmemente ao \textbf{método} científico.
Sendo assim, foi considerado \textbf{um} dos fundadores da Química Moderna.
O seu principal foco de trabalho foram os gases, sua obra mais famosa, o livro ; THE SCEPTICAL CHYMIST ( O químico cético ).
Desta forma, mudou a interpretação da Química de seu tempo, além de criticar o trabalho dos alquimistas, onde a ciência passou a ser chamada de Química.
No \textbf{século} XVIII, o \textbf{século} do Iluminismo, a Química tem sua certidão de nascimento, na era da revolução industrial, capitalismo em pleno desenvolvimento, observa -se a expansão da ciência tecnológica.
Alguns historiadores conclamam Antonie Laurent Lavoisier como o “ Pai da Química Moderna ”.
Finalmente no \textbf{século} XIX, a Química se consolidou com a proposição de \textbf{um} modelo por Dalton \textbf{que} \textbf{influenciou} o desenvolvimento posterior da química. ( GREENBERG, 2009 ).
De acordo com as Diretrizes Curriculares Nacionais da Educação Básica ( DCNEM, 2013 ), foi no final do \textbf{século} XIX, com o surgimento dos laboratórios de pesquisa, \textbf{que} a Química se consolidou como a principal \textbf{disciplina} associada aos efetivos resultados na indústria.
No \textbf{século} XX, a Química e todas as outras Ciências Naturais tiveram \textbf{um} grande desenvolvimento, em especial nos Estados Unidos, Inglaterra e Alemanha.
Esses países destacaram -se no desenvolvimento da Ciência, no intuito de estabelecer e posteriormente, manter a \textbf{influência} científica \textbf{que} pudesse garantir diferentes formas de poder e controle bélico mundial, essenciais nas tensões \textbf{vividas} no \textbf{século} XX.
Diversos foram os investimentos desses países em áreas como : obtenção de medicamentos, indústria bélica, estudos nucleares, estrutura atômica e formação das moléculas, mecânica quântica, dentre outras \textbf{que} estreitaram as relações entre a ciência e a indústria. ( GREENBER, 2009 ).
Podemos elucidar \textbf{que} é papel do professor abordar de maneira diferente do ensino tradicional o conceito da História da Química e suas implicações no Ensino Médio.
Destaca -se a importância de inovar, impulsionar novos modelos e estratégias de trabalho em sala de aula, buscando metodologias ativas com foco no protagonismo e o projeto de vida do estudante, estruturado nos quatros pilares da educação, \textbf{que} são : aprender a conhecer, aprender a fazer, aprender a ser e aprender a conviver.
A potencialidade de ensinar a partir da perspectiva da formação de \textbf{um} \textbf{sujeito} capaz de tomar decisões é importante para a constituição de \textbf{uma} educação emancipadora e o ensino de Química não pode se isentar dessa responsabilidade.
Dessa forma, deve estar presente nas aulas de Química a decisão de trabalhar as propostas a partir de temáticas \textbf{que} possibilitem aos estudantes se posicionarem de maneira consciente frente aos fatos do cotidiano.
A lei de Diretrizes e Bases da Educação Brasileira ( LDB 9.394/96 ) e os Parâmetros Curriculares Nacionais ( PCNEM, 1999 ) \textbf{salientam} \textbf{que} o enfoque tecnicista da aprendizagem, aparece na segunda metade do \textbf{século} XX, no Brasil em 1960-1970 – \textbf{que} \textbf{implica} em priorizar a transmissão e recepção de informações - ainda \textbf{hoje} paradigmas da educação no Brasil, não satisfaz mais às necessidades da sociedade brasileira.
Assim, os PCNEM enfatizam \textbf{que} o ensino, nos níveis Fundamental e Médio, deve favorecer a formação geral do indivíduo, como descritos no art. 5 e 6 das Diretrizes Curriculares Nacionais Gerais para o Ensino Médio ( DCNEM, 2018 ), ao invés de se voltar para a formação de \textbf{um} profissional específico.
Segundo a concepção dos processos de ensino aprendizagem sugerida nesses documentos, o aprender \textbf{implica} no desenvolvimento das capacidades de pesquisa e criação, não \textbf{baseado} apenas no exercício da \textbf{memorização}.
Neste contexto, a Base Nacional Curricular Comum ( BNCC, 2018 ), em conformidade com a alteração da LDB 9.394/96 por força da Lei n{\EmojiFont °} 13.415/2017, busca de forma intencional a formação integral do estudante nos aspectos físicos, cognitivos e sociemocionais por meio de processos educativos significativos \textbf{que} promovam a autonomia, o comportamento cidadão e o protagonismo na construção de seu projeto de vida.
Diante dessa perspectiva, o uso de temas Químicos-Sociais, tais como : Química ambiental ; Metais, metalurgia e galvanoplastia ; Química dos materiais sintéticos ; Recursos energéticos ; Alimentos e aditivos químicos ; Minerais ; Energia nuclear ; Medicamentos ; Química na agricultura ; Bioquímica ; Água ; Processos industriais ; Petróleo petroquímica ; Drogas ; Sabões e detergentes ; Plásticos ; Tintas ; Geoquímica ; Vestuários ; Materiais importados pelo Brasil ; Química da arte e Recursos naturais, no ensino de Química se torna \textbf{uma} ferramenta potente para essa \textbf{tarefa}, pois os temas são extraídos das relações do \textbf{sujeito} com seu contexto ( local, regional, nacional ou mundial ) e seu desenvolvimento, e \textbf{desempenham} papel de fundamental importância no ensino de química para formar o cidadão, pois propiciam a \textbf{contextualização} do conteúdo químico com o cotidiano do estudante. ( Santana \textbf{et} al, 2014 ).
Além disso, permitem o desenvolvimento das competências e habilidades básicas relativas à cidadania, como a participação e a capacidade de tomada de decisão, pois trazem para sala de aula discussões de aspectos sociais relevantes, \textbf{que} \textbf{exigem} dos estudantes \textbf{um} \textbf{posicionamento} crítico quanto a sua solução. ( Santos \& Schnetzler, 1996, p. 30 ).
Neste sentido, Caro(a) professor(a), o currículo de Sergipe da área Ciências da Natureza e suas Tecnologias, na componente curricular de Química foi estruturado de forma \textbf{simples} e clara, para os senhores \textbf{que} de fato irão colocar “ a mão na massa ”.
O grande desafio é a mobilização das competências gerais e específicas e habilidades por meio dos objetos do conhecimento, \textbf{que} estão distribuídos nos quatro primeiros semestres considerando os ensinos adquiridos no fundamental e \textbf{que} serão aprofundados no ensino médio.
O organizador curricular de Ciências da Natureza apresenta sete colunas, representadas por : competências, habilidades, semestres, unidades temáticas e as três componentes curriculares ( Biologia, Química e Física ).
Das 26 habilidades de Ciências da Natureza presente na BNCC de 2018, química contemplou 25 habilidades.
Ao mobilizá -las, observem os níveis cognitivos de complexidade da taxonomia de Bloom, analise se o estudante apresenta habilidades do pensamento de ordem inferior para \textbf{que} possa desenvolver a habilidade de pensamento de ordem superior, correspondente à habilidade do currículo \textbf{que} será mobilizada.
Caso o estudante não apresente, utilizem sequências didáticas, oficinas temáticas e experimentação como estratégicas \textbf{metodológicas}, \textbf{focadas} no desenvolvimento do protagonismo e projeto de vida, assegurando assim \textbf{uma} educação integral do estudante.
Os objetos do conhecimento da componente curricular de Química foram distribuídos buscando almejar ao final do Ensino Médio \textbf{que} o estudante \textbf{consiga} aprender a estrutura da matéria e sua transformação \textbf{que} deve ser trabalhado no primeiro e segundo semestre, compreender os aspectos dinâmicos da matéria por meio do estudo da cinética química, equilíbrio químico, eletroquímica \textbf{que} será \textbf{visto} no terceiro semestre e no quarto semestre aprender sobre a química do carbono e sobretudo mobilizar as habilidades, valores e atitudes.
Desta forma, esperamos formar jovens autônomos com pensamento científico, crítico e criativo.

% matched lemmas: adiante, basear, conhecido, conseguir, contextualização, desempenhar, disciplina, et, exigir, focar, hoje, homem, humanidade, implicar, influenciar, influência, memorização, metodológico, moderno, método, permear, posicionamento, que, salientar, simples, sujeito, século, sólido, tarefa, um, ver, viver
\end{textsample}
